% ===== 原文 PDF pp.59--63（印刷 pp.58--62）=====
% tr-05（3/3）。首行紧接 tr05b 末句（同一自然段，文件间不留空行）。
% [原文 p.58]
当穿刺点为正则时，对 $A_N$ 型 Class S 理论，每个穿刺点的极点结构都对应一张杨图\footnote{在这种情况下，按 [36] 引入的记号，该穿刺点也被称为「正则」。需要强调，这两种「正则穿刺点」的概念并不重合；在本文中，我们所说的「正则」穿刺点，是指 Hitchin 场在该点处行为如式 (4.12) 的那些穿刺点。}，它编码整数 $N$ 的一个分拆。于是整体对称性可以直接从该杨图读出。

考虑式 (4.9) 中的 SW 曲线。我们可以用一组正整数 $\{p_j\}=\{p_2,\ldots,p_N\}$ 来定义穿刺点的 \textbf{极点结构}。它们是 $\phi_j$ 在穿刺点处可以容许的极点阶数。由极点的结构，我们如下找出与穿刺点相联系的味对称群 [33,36]。

首先，对 A 型穿刺点，我们把一张杨图与极点结构 $\{p_k\}$ 相联系：
\begin{enumerate}
  \item 从第一行有两个方格的一张杨图出发。
  \item 对每个 $k=3,\ldots,N$：
  \begin{itemize}
    \item 若 $p_k-p_{k-1}=1$，在当前行上加一个方格；
    \item 若 $p_k-p_{k-1}=0$，在其下方另起一行，放一个方格。
  \end{itemize}
\end{enumerate}
于是整体味对称群为
\begin{equation}
  G = \mathrm{S}\left(\prod_h \mathrm{U}(n_h)\right),
  \tag{4.16}
\end{equation}
其中 $n_h$ 是与该穿刺点相联系的杨图中高度为 $h$ 的列数。\footnote{整体味对称性可能会增强。为了找出正确的整体对称群，例如计算该理论的超共形指标或 Hilbert 级数是有帮助的。}
% [Tutor Note: (4.16) 的结构：$\prod_h\mathrm{U}(n_h)$ 是分拆 $\{h^{n_h}\}$ 对应 Levi 子群的味对称群，前面的 $\mathrm{S}(\cdot)$ 是去掉整体 $\mathrm{U}(1)$（即行列式为 1 的约束），与 $T_2$ 例子里 $\mathrm{S}(\mathrm{U}(1)\times\mathrm{U}(2)\times\mathrm{U}(1))$ 的写法一致。例 4.2 的杨图有 1 个高三的列、2 个高二的列、1 个高一的列，故 $n_1=1,\,n_2=2,\,n_3=1$，正是 $\mathrm{S}(\mathrm{U}(1)\times\mathrm{U}(2)\times\mathrm{U}(1))$。]

\begin{examplebox}{例 4.2：}
  作为例子，考虑 $\{p_k\}=\{1,2,3,3,4,5,5\}$。我们于是从一张如下的杨图出发：\framebox[1.3em]{\phantom{0}}\framebox[1.3em]{\phantom{0}}。由于 $p_3-p_2=1$，我们在同一行上加一个方格；对 $p_4-p_3=1$ 也照做，但对 $p_5-p_4$，我们另起一行。完成该程序后，我们得到如下杨图：\begin{tabular}[b]{@{}c@{}c@{}c@{}c@{}}\framebox[1.3em]{0} & \framebox[1.3em]{1} & \framebox[1.3em]{2} & \framebox[1.3em]{3}\\\framebox[1.3em]{3} & \framebox[1.3em]{4} & \framebox[1.3em]{5}\\\framebox[1.3em]{5}\end{tabular}。其中有 1 个高度为三的列、2 个高度为二的列和 1 个高度为一的列。因此整体对称性是 $\mathrm{S}\left(\mathrm{U}(1)\times\mathrm{U}(2)\times\mathrm{U}(1)\right)$。
\end{examplebox}

% [Tutor Note: 例 4.2 里按算法的记账：$p_1=0$ 不列入集合，故第一行先放两个方格对应 $p_1=0\to p_2=1$；此后 $p_3-p_2=1$、$p_4-p_3=1$ 各加一格（第一行 4 格），$p_5-p_4=3-3=0$ 另起一行，$p_6-p_5=1$、$p_7-p_6=1$ 各加一格（第二行 3 格），$p_8-p_7=5-5=0$ 再另起一行（第三行 1 格），共 8 格 = $N=8$。]

上述程序也可以反过来做，由杨图重建各微分 $\phi_k$ 的极点结构 $\{p_k\}$（例如见 [36]）：
% [原文 p.59]
\begin{enumerate}
  \item 从最长的一行开始，用整数标记杨图的 $N$ 个方格，并把该行最左边的方格标为 0。
  \item 沿某一行向右移动时，把标号逐个加一。
  \item 移到它下面的一行，给它最左边的方格标上与上一行最右方格相同的标号；在该行重复上述标记程序。
  \item 所得的 $N$ 个标号序列即为 $\{p_1,\ldots,p_N\}$，其中 $p_1=0$，而 $p_2$ 必然等于 1。
\end{enumerate}
此外，若一个正则穿刺点对所有 $k$ 都满足 $p_k=k-1$，则称它为 \textbf{极大}或 \textbf{完全}穿刺点，其杨图是一行水平的方格。$\SU(N)$ 的杨图最多只能有 $N-1$ 行，与之相联系的穿刺点称为 \textbf{极小}或 \textbf{简单}穿刺点。

\begin{examplebox}{例 4.3：}
  作为例子，我们可以从例 4.2 的杨图出发，但不给方格标号：\begin{tabular}[b]{@{}c@{}c@{}c@{}c@{}}\framebox[1.3em]{\phantom{0}} & \framebox[1.3em]{\phantom{0}} & \framebox[1.3em]{\phantom{0}} & \framebox[1.3em]{\phantom{0}}\\\framebox[1.3em]{\phantom{0}} & \framebox[1.3em]{\phantom{0}} & \framebox[1.3em]{\phantom{0}}\\\framebox[1.3em]{\phantom{0}}\end{tabular}。我们按刚才给出的规则标记该杨图，结果（显然）是下面这张杨图：\begin{tabular}[b]{@{}c@{}c@{}c@{}c@{}}\framebox[1.3em]{0} & \framebox[1.3em]{1} & \framebox[1.3em]{2} & \framebox[1.3em]{3}\\\framebox[1.3em]{3} & \framebox[1.3em]{4} & \framebox[1.3em]{5}\\\framebox[1.3em]{5}\end{tabular}。极点结构为 $\{p_k\}=\{1,2,3,3,4,5,5\}$。这可以与式 (4.4) 中 Hitchin 场的显式形式相联系。该杨图的行与列分别为 $s_i=\{4,3,1\}$ 与 $t_i=\{3,2,2,1\}$。

  当超多重态无质量时，穿刺点处的 Hitchin 场 $\Phi_z$ 具有如下形式的留数（另见 [62]）
  \begin{equation}
    \Res\Phi_z \sim J_{s_1}\oplus J_{s_2}\oplus J_{s_3}
    = \begin{pmatrix}
      J_{s_1} & 0 & 0\\
      0 & J_{s_2} & 0\\
      0 & 0 & J_{s_3}
    \end{pmatrix},
    \tag{4.17}
  \end{equation}
  其中
  \begin{equation}
    J_4 =
    \begin{pmatrix}
      0 & 1 & 0 & 0\\
      0 & 0 & 1 & 0\\
      0 & 0 & 0 & 1\\
      0 & 0 & 0 & 0
    \end{pmatrix},\qquad
    J_3 =
    \begin{pmatrix}
      0 & 1 & 0\\
      0 & 0 & 1\\
      0 & 0 & 0
    \end{pmatrix},\qquad
    J_1 = 0\,,
    \tag{4.18}
  \end{equation}
  一般地，$J_s$ 是 $s\times s$ 的 Jordan 块矩阵。于是很容易证明：把式 (4.17) 代入式 (4.4) 并与式 (4.9) 比较，$\phi_k(z)$ 的极点阶数
% [原文 p.60]
  恰好由 $\{p_k\}$ 给出。

  在超多重态有质量的情形，这个矩阵值 1-形式的留数应共轭于矩阵
  \begin{equation}
    \diag\left(m_1,m_1,m_1,m_2,m_2,m_3,m_3,m_4\right),
    \tag{4.19}
  \end{equation}
  并受无迹条件的约束。更一般地，对列高为 $\{t_i\}$ 的杨图，Hitchin 场的留数共轭于如下对角矩阵
  \begin{equation}
    \diag\left(\underbrace{m_1,\ldots,m_1}_{t_1},
    \underbrace{m_2,\ldots,m_2}_{t_2},\ldots\right),
    \tag{4.20}
  \end{equation}
  其中各质量满足
  \begin{equation}
    \sum_i t_i m_i = 0\,.
    \tag{4.21}
  \end{equation}
  这样的留数可以与式 (4.16) 中味对称性的质量参数相联系。就本例而言，味对称性是
  \begin{equation}
    \mathrm{S}\left(\mathrm{U}(1)\times\mathrm{U}(2)\times\mathrm{U}(1)\right).
    \tag{4.22}
  \end{equation}
  于是有一个质量参数 $m_1$ 与第一个 $\mathrm{U}(1)$ 相联系，两个质量参数 $m_2$ 与 $m_3$ 属于 $\mathrm{U}(2)$，还有一个质量参数 $m_4$ 属于第二个 $\mathrm{U}(1)$，它们之间有一个线性关系，与 (4.19) 一致。推广到任意穿刺点是直截了当的。
\end{examplebox}

% [Tutor Note: (4.17) 中 $s_i$ 是杨图各行长度（$s=\{4,3,1\}$），故 Jordan 块取 $J_4\oplus J_3\oplus J_1$；$t_i$ 是各列高度（$t=\{3,2,2,1\}$），故有质量的情形下对角元按列高重复，且 $\sum_i t_im_i=\Tr\,\diag(\cdots)=0$ 正是无迹条件。两种写法描述的是同一个余伴随轨道：半单（有质量）时由列高给出本征值的重数，幂零（无质量）时由行长方体的 Jordan 块给出。]

\begin{exercisebox}{习题 4.2}
  证明：对例 4.3 把式 (4.17) 代入式 (4.4) 并与式 (4.9) 比较，$\phi_k(z)$ 的极点阶数恰好由 $\{p_k\}$ 给出。
\end{exercisebox}

一如既往，若想对穿刺点及其与低能理论整体对称性的关系获得一些直觉，从弦论出发思考是有启发的。正如 4.3 节所解释的，穿刺点在 M 理论中起源于在特殊点 $p_i$ 处与 Gaiotto 曲线横截的缺陷 M5 膜。约化到 IIA 型弦论后 [28]，这些膜变成半无限的 D4 膜或 NS5 膜。一叠 $N$ 张 D4 膜上有低能 $\SU(N)$ 规范理论。事实上，规范群本应是 $\mathrm{U}(N)$，只是其中一个 $\mathrm{U}(1)$ 因子——它对应于这叠 D4 膜作为整体沿 NS5 膜竖直方向的运动（同时在其左侧与右侧）——与谱退耦（这一点稍后会变得相关）。

% [原文 p.61]
D4 膜是半无限的，这使得这些规范场的动能项，与悬挂在两张 NS5 膜之间的有限 D4 膜所对应的同类项相比，可以忽略。因此，一叠 $N$ 张半无限 D4 膜会在低能理论中给出一个 $\SU(N)$ 味对称性。IIA 型的 NS5 膜没有这样的低能规范理论描述，因为单张 NS5 膜的世界体动力学由二形式 $B_{\mu\nu}$ 及其超对称伙伴支配，而让若干张 NS5 膜重合，会在张量分支的原点处生成一个含无张力弦的 6d 理论（它在低能极限下定义一个 $\cN=(2,0)$ SCFT）。不过，紧化到 4 维之后，$B$ 场的约化会为每张五维膜给出一个涌现的 $\mathrm{U}(1)$ 规范场。在我们的情形中，NS5 膜是无限的但并不重合，这意味着它们每一张都给出一个 $\mathrm{U}(1)$ 整体对称性。总而言之，我们得出结论：Class S 的完全穿刺点与简单穿刺点在 IIA 型膜图中都有简单的解释，分别对应于左右两端的半无限 D4 膜与竖直的 NS5 膜。

\begin{center}
  \includegraphics[width=0.9\linewidth]{fig9.png}
\end{center}
\figtitle{9}{Class S 中正则穿刺点与 IIA 型膜图之间的对照表。左图画出的是含 D4--NS5--D6 膜的 HW 装置最右端部分，右图画的是黎曼面上相应的穿刺点。NS5 膜及其相应的极小穿刺点用蓝色画成。}

% [Tutor Note: 读图 9 的要点：左侧每条水平线是一张 D4 膜，右端的圆圈叉号是 D6 膜；蓝色竖线是 NS5 膜；左侧的竖排省略号表示向膜图内部无限延伸。右侧黎曼面上的杨图（蓝色）对应蓝色 NS5 膜给出的极小穿刺点，黑色杨图（4/3/1 形状）对应由 D6 膜标出的更复杂的正则穿刺点；每一列的高度等于该处终结的 D4 膜数。]
其他与更复杂杨图相联系的正则穿刺点，只要我们在图像中加入 D6 膜，也很容易在 IIA 型中理解。做法是：对给定的穿刺点，先加入与标记该穿刺点的杨图中的列数一样多的 D6 膜，然后让 D4 膜终结在 D6 膜上，并按每一列方格数所指示的方式分配它们（例子见图 9）。现在既然加入了 D6 膜，读取 HW 装置整体对称性的方式就与先前只有半无限 D4 膜的情形不同了。这是因为悬挂在 NS5 膜与 D6 膜之间的 D4 膜没有低能
% [原文 p.62]
规范理论自由度；事实上，它的所有无质量自由度都被 NS5 膜与 D6 膜所施加的边界条件冻结了。相反，在整体对称性下带荷的场来自悬挂在 D4 膜与 D6 膜之间的基本弦，而这些整体对称性的背景规范场来自 D6 膜的世界体。方便的做法是作若干次 Hanany--Witten 移动，把 D6 膜移到膜图的内部，并使「冻结的」D4 膜消失。在图 9 的例子中，我们有两张各带一张 D4 膜的 D6 膜：我们可以把它们移到 NS5 膜之间的第一个区间并湮灭这两张 D4 膜；两张重合的 D6 膜给出一个 $\mathrm{U}(2)$ 整体对称性因子。同样，对于有两张（相应地，三张）D4 膜终结其上的那张 D6 膜，我们需要作两次（相应地，三次）Hanany--Witten 移动，把它带到第二（相应地，第三）个 NS5 膜区间：两种情形下都给出 $\mathrm{U}(1)$ 整体对称性因子。注意，只要我们按上一段所解释的去掉整体的 $\mathrm{U}(1)$ 因子，这恰好就是我们在 (4.16) 中得到的整体对称性。

\begin{center}
  \includegraphics[width=0.5\linewidth]{fig10.png}
\end{center}
\figtitle{10}{带 4 个穿刺点的黎曼球面。}

现在让我们回到 6 维，考虑包含比各穿刺点质量更多参数的更有趣的理论。一种做法是给球面再加一个穿刺点：$\cC$ 上的全纯坐标重定义只允许我们把其中三个穿刺点的位置固定到 $z=0,1,\infty$。第四个穿刺点的位置由 $q$ 给出，它参数化带穿刺点球面的复结构，是一个可以发生有趣物理的新参数。特别地，$q$ 是无量纲的，它定义了 4d 理论的一个恰当边缘参数。更一般地，UV 曲线 $\cC$ 的复结构模，在 Class S 构造中为所得的 4d $\cN=2$ SCFT 定义一个共形（子）流形。
% [Tutor Note: 这里的「边缘」是指 $q$ 不进入任何 beta 函数：$\cC$ 的复结构模在 4d 中对应超势的恰当边缘耦合（对 $\cN=2$ 就是规范耦合 $\tau$ 的实部/虚部的组合），而 Class S 中它恰是 $q$；因此 4d 共形流形的一部分就是带穿刺点黎曼面的复结构模空间（模去 $\mathrm{PSL}(2,\C)$ 后剩下的 $n-3$ 个模）。]

更具体地说，让我们考虑紧化在图 10 中球面上的 $\fsu(2)$ 6d $\cN=(2,0)$ 理论，其中 4 个穿刺点中的每一个都是 Hitchin 场的极点，如式 (4.12) 所示。
% [边界] 本句已写完（原文 p.62 末行 "as in equation" 起、p.63 首行 "(4.12)." 止）。下一位 tr-06 请从 PDF p.64 首行之后（即 "Exercise 4.3"）开始，勿重写 "(4.12)." 与本句。

